\documentclass[12pt]{article}
\usepackage{graphicx}
\usepackage[margin=2cm, a4paper]{geometry}
\usepackage{setspace}
\usepackage{pdfpages}
\usepackage{float}
\usepackage{amsmath}
\usepackage{fancyvrb}
\usepackage{amssymb}
\usepackage{minted}
\usepackage{enumitem}
\usepackage[colorlinks,linkcolor=blue]{hyperref}


\renewcommand{\contentsname}{Contents}
\renewcommand{\figurename}{Figure}
\renewcommand{\tablename}{Table}
\hypersetup{
    colorlinks=true,
    linkcolor=black,
    filecolor=magenta,      
    urlcolor=blue,
}
\newcommand{\mytitle}{Abstract Algebra II - Homework 1}
\newcommand{\myauthor}{B13902022 Yu-Chi,Lai}

\usepackage{fancyhdr}
\pagestyle{fancy}
\fancyhead{}
\fancyhead[L]{\mytitle}
\fancyhead[R]{\myauthor}

\title{\mytitle}
\author{\textbf{\myauthor}}
\date{}

\begin{document}

\maketitle

\section{}
Consider the evaluation map defined by substituting $x$ with $u$:
$$\varphi: K[x]\rightarrow L,\ g(x)\mapsto g(u)$$
The image of this homomorphism, $\mathrm{Im}{(\varphi)}$, is $K[u]$ by definition. Because $u$ is algebraic over $K$, there exists at least one non-zero polynomial in $K[x]$ having $u$ as a root. Hence, $\ker{(\varphi)}\ne \{0\}$.

Since $K$ is a field, thus $K[x]$ is a PID. Let the monic generator of $\ker{(\varphi)}$ be $f(x)$ (We can do so because multiplying $f$ by a constant, the ideal generated is unchanged), i.e., $\ker{(\varphi)}=(f(x))$. By the first isomorphism for rings, we have:
$$K[x]/(f(x))\cong K[u]$$
\subsection{(b)}
Since $K[u]$ is a subring of $L$, and $L$ is an integral domain, so $K[u]$ is also a integral domain. Since $K[x]/(f(x))$ is an integral domain by the isomorphism, $(f(x))$ is a prime ideal, and every non-zero prime ideal in a PID is maximal.

Because $(f(x))$ is maximal, its generator $f(x)$ is irreducible, and furthermore, $K[x]/(f(x))$ is a field.

Finally, $f(x)$ satisfies two properties. Since $f\in \ker{\varphi}$, we have $f(u)=0$. And for any $g\in K[x]$ such that $g(u)=0$, it means that $g\in \ker{\varphi}=(f(x))$, hence $g(x)$ is divisible by $f(x)$. The converse is also true, if $g(x)$ is divisible by $f(x)$, it means $g(x)\in (f(x))=\ker{(\varphi)}$, so $g(u)=0$.

\subsection{(a)}
Since $K(u)$ contains all element of the forms $\frac{g(u)}{1}$ where $g(x)\in K[x]$, $K[u]$ automatically lies in $K(u)$, so we have $K[u]\subseteq K(u)$.

But from (b) we know $K[u]$ is a field containing both $K$ and $u$, and $K(u)$ is the smallest subfield of $L$ meets this requirement. Hence, $K[u]=K(u)$, thus $K[u]\cong K(u)$.
\section{}
Since $u$ is algebraic over $K$, then there exists a unique monic irreducible polynomial $f\in K[x]$ such that $f(u)=0$. Let the isomorphism from $K(u)$ to $K(u')$ be $\varphi$, let $f(x)=a_0+a_1x+a_2x^2+\ldots+ a_nx^n$.

$$f(u)=a_0+a_1u+a_2u^2+\ldots+a_nu^n\in K(u)$$
If we map $f(u)$ to $K(u')$ using $\varphi$, we get:
\begin{align*}    
\varphi(f(u))&=\varphi(a_0+a_1u+a_2u^2+\ldots+a_nx^n) \\
&=\varphi(a_0)+\varphi(a_1)\varphi(u)+\varphi(a_2)\varphi(u^2)+\ldots+\varphi(a_n)\varphi(u^n) \\
&=a_0+a_1u'+a_2u'^2+\ldots+a_nu'^n \\
&=f(u')=0
\end{align*}
Hence, $u$ and $u'$ are roots of the same irreducible polynomial $f\in K[x]$.

\end{document}
% how to display codes?
% \begin{minted}[frame=lines,framesep=2mm,baselinestretch=1.2,linenos,breaklines]{python}
% \end{minted}

% how to display images?
% \begin{figure}[H]
%     \centering
%     \includegraphics[width=0.5\linewidth]{}
%     \caption{Caption}
% \end{figure}
% test